\documentclass[11pt, a4paper]{article}

% --- Paquetes Fundamentales ---
\usepackage[utf8]{inputenc}
\usepackage[T1]{fontenc}
\usepackage[spanish,es-tabla]{babel}
\usepackage[margin=2.5cm, headheight=25pt]{geometry} % Corregido headheight
\usepackage{graphicx}
\usepackage{fancyhdr}
\usepackage{xcolor}
\usepackage{colortbl} % Añadido para el comando \cellcolor
\usepackage{microtype}
\usepackage[colorlinks=true, linkcolor=ucbblue, urlcolor=ucbblue]{hyperref}

% --- Matemáticas y Electrónica ---
\usepackage{amsmath, amssymb}
\usepackage{siunitx}
\usepackage[american, RPvoltages]{circuitikz}

% --- Elementos de Diseño Pedagógico y Código ---
\usepackage{tcolorbox}
\usepackage{enumitem}
\usepackage{listings}
\usepackage{tabularx}

% --- Configuración de Colores y siunitx ---
\definecolor{ucbblue}{RGB}{0, 51, 102}
\definecolor{codegreen}{rgb}{0,0.6,0}
\definecolor{codepurple}{rgb}{0.58,0,0.82}
\sisetup{
    output-decimal-marker = {,},
    per-mode = symbol,
    inter-unit-product = \ensuremath{{}\cdot{}}
}

% --- Macros y Estilos ---
\tcbset{
    caja_fase/.style={colback=blue!3, colframe=ucbblue, fonttitle=\bfseries, boxrule=1pt, arc=3pt},
    caja_alerta/.style={colback=orange!5, colframe=orange!80!black, fonttitle=\bfseries, boxrule=0.8pt, arc=3pt}
}

\lstset{
    language=Python,
    basicstyle=\ttfamily\footnotesize,
    keywordstyle=\color{blue}\bfseries,
    commentstyle=\color{codegreen}\itshape,
    stringstyle=\color{codepurple},
    numbers=left, numberstyle=\tiny\color{gray},
    frame=single, backgroundcolor=\color{gray!5},
    breaklines=true, showstringspaces=false
}

% --- Estilo de Página ---
\pagestyle{fancy}
\fancyhf{}
\lhead{\textcolor{ucbblue}{\textbf{IMT-121 | Guía de Clase y Material Maestro}}}
\rhead{\textbf{Sesión 08 -- Semana 4}}
\cfoot{Página \thepage}
\renewcommand{\headrulewidth}{0.4pt}
\renewcommand{\headrule}{\hbox to\headwidth{\color{ucbblue}\leaders\hrule height \headrulewidth\hfill}}

\begin{document}

\begin{center}
    {\LARGE \textbf{GUION DE CLASE Y MATERIAL MAESTRO}} \\[0.2cm]
    {\Large \textbf{SESIÓN 08}} \\[0.4cm]
    \normalsize
    \textbf{Asignatura:} IMT-121 Circuitos Electrónicos I \\
    \textbf{Tema:} Estrategia Dual de Thévenin: Del Sensor Pasivo al Acondicionamiento Activo ($V_{\text{test}} / I_{\text{test}}$) \\
    \textbf{Gestión:} Semana 4 — Sesión 1 (90 Minutos) \\
    Universidad Católica Boliviana ``San Pablo'' — Sede Tarija \\[0.2cm]
    \textbf{Docente:} M.Sc. Ing. Bernardo Quiroga Turdera
    \vspace{0.4cm}
\end{center}
\hrule
\vspace{0.5cm}

\section*{Estructura Pedagógica de la Sesión (90 Minutos)}
\begin{center}
    \renewcommand{\arraystretch}{1.3}
    \begin{tabularx}{\textwidth}{|X|}
        \hline
        \cellcolor{ucbblue}\textcolor{white}{\textbf{ESTRUCTURA DE CLASE EN DOS FASES INTEGRADAS}} \\
        \hline
        \cellcolor{gray!10}\textbf{FASE 1: DOMINIO PASIVO (00–45 min) — El Sensor Desbalanceado} \\
        $\bullet$ \textbf{00–10 min:} Problematización: Modelado Thévenin en sensores de deformación. \\
        $\bullet$ \textbf{10–25 min:} Deducción analítica de $V_{oc}$ y $R_{th}$ pasivo por división y reducción. \\
        $\bullet$ \textbf{25–40 min:} Condición de Máxima Transferencia de Potencia y Eficiencia del $50\%$. \\
        $\bullet$ \textbf{40–45 min:} Verificación en Python del barrido de carga ($R_L$ vs $P_L$). \\
        \hline
        \cellcolor{gray!10}\textbf{FASE 2: EXTENSIÓN ACTIVA (45–90 min) — Preamplificación y Caja Negra} \\
        $\bullet$ \textbf{45–55 min:} Inserción de etapa activa (VCVS): ¿Por qué colapsa apagar fuentes? \\
        $\bullet$ \textbf{55–70 min:} Algoritmo de la Fuente de Prueba Externa ($V_{\text{test}} / I_{\text{test}}$) en pizarra. \\
        $\bullet$ \textbf{70–80 min:} Comparación cuantitativa: $R_{th}$ pasiva vs. $R_{th}$ activa. \\
        $\bullet$ \textbf{80–90 min:} Taller LTspice (\texttt{.tf} y \texttt{.op}) + Dinámica en Tríadas y Cierre. \\
        \hline
    \end{tabularx}
\end{center}

\vspace{0.3cm}
\begin{tcolorbox}[caja_fase, title=FASE 1: DOMINIO PASIVO (00 A 45 MIN)]
\textbf{Modelado Thévenin en Puentes de Medición y Máxima Transferencia}
\end{tcolorbox}

\textbf{Contexto Físico:} \\
Una celda de carga industrial o un sensor de bioimpedancia dérmica se modela mediante un Puente de Wheatstone. Cuando el sensor detecta una variable física, su resistencia cambia en $\Delta R$, generando un voltaje diferencial que alimenta la siguiente etapa.

\begin{center}
    \begin{circuitikz}[scale=0.9, transform shape]
        \draw (0,2) node[left]{Terminal A ($V_A$)} to[R, l=$R_3 {=} \qty{350}{\ohm}$] (2,0) node[below]{MASA / GND (\qty{0}{\volt})};
        \draw (0,2) to[R, l=$R_1 {=} \qty{350}{\ohm}$] (2,4) node[above]{NODO DE ALIMENTACIÓN ($V_s {=} \qty{10}{\volt}$)};
        \draw (2,4) to[R, l=$R_2 {=} \qty{350}{\ohm}$] (4,2) node[right]{Terminal B ($V_B$)};
        \draw (4,2) to[R, l=$R_4 {=} \qty{360}{\ohm} \, (R{+}\Delta R)$] (2,0);
        \draw (2,0) node[ground]{};
        \draw (2,4) to[short, *-] (2,4.2);
        \draw (0,2) to[short, *-o] (-0.5,2);
        \draw (4,2) to[short, *-o] (4.5,2);
    \end{circuitikz}
\end{center}

\subsection*{1. Deducción Analítica de $V_{th}$ (Voltaje de Circuito Abierto)}
El voltaje de Thévenin entre los terminales $A$ y $B$ corresponde a la diferencia de potencial generada por los dos divisores de tensión independientes:
\begin{align*}
    V_A &= \qty{10}{\volt} \cdot \left( \frac{350\,\Omega}{350\,\Omega + 350\,\Omega} \right) = 10 \cdot \frac{350}{700} = \mathbf{\qty{5.000}{\volt}} \\
    V_B &= \qty{10}{\volt} \cdot \left( \frac{360\,\Omega}{350\,\Omega + 360\,\Omega} \right) = 10 \cdot \frac{360}{710} = \mathbf{\qty{5.0704}{\volt}}
\end{align*}
El voltaje en circuito abierto visto desde los terminales $B-A$ es:
\begin{equation}
    V_{th} = V_{oc} = V_B - V_A = \qty{5.0704}{\volt} - \qty{5.0000}{\volt} = \mathbf{\qty{0.07042}{\volt} = \qty{70.42}{\milli\volt}}
\end{equation}

\subsection*{2. Deducción Analítica de $R_{th}$ Pasiva (Apagando Fuentes)}
Para calcular la resistencia equivalente vista entre los terminales $A$ y $B$:
\begin{itemize}
    \item Apagamos la fuente independiente de tensión: $V_s = \qty{0}{\volt} \implies$ Cortocircuito a Tierra ($GND$).
    \item La parte superior del puente queda unida al riel de Tierra.
    \item Vista desde los terminales $A$ y $B$, la red se desacopla en dos ramas en paralelo conectadas en serie entre sí:
\end{itemize}

\begin{center}
    \begin{circuitikz}[scale=0.9, transform shape]
        \draw (0,0) node[left]{Terminal A} to[short, o-] (1,0) -- (1,1) to[R, l=$R_1$] (3,1) -- (3,0);
        \draw (1,0) -- (1,-1) to[R, l=$R_3$] (3,-1) -- (3,0);
        \draw (3,0) node[ground]{};
        \draw (3,0) -- (4,1) to[R, l=$R_2$] (6,1) -- (6,0);
        \draw (3,0) -- (4,-1) to[R, l=$R_4$] (6,-1) -- (6,0);
        \draw (6,0) to[short, -o] (7,0) node[right]{Terminal B};
    \end{circuitikz}
\end{center}

\begin{align*}
    R_A &= R_1 \parallel R_3 = \frac{350 \cdot 350}{350 + 350} = \mathbf{175.0\,\Omega} \\
    R_B &= R_2 \parallel R_4 = \frac{350 \cdot 360}{350 + 360} \approx \mathbf{177.465\,\Omega} \\
    R_{th} &= R_A + R_B = 175.0\,\Omega + 177.465\,\Omega = \mathbf{352.465\,\Omega}
\end{align*}

\subsection*{3. Teorema de Máxima Transferencia de Potencia}
Para transferir la máxima potencia pasiva del sensor a una carga $R_L$:
\begin{equation}
    \mathbf{R_L = R_{th} = 352.465\,\Omega}
\end{equation}
Voltaje en la carga acoplada:
\begin{equation}
    V_L = V_{th} \cdot \left( \frac{R_L}{R_{th} + R_L} \right) = \frac{V_{th}}{2} = \mathbf{35.21\text{ mV}}
\end{equation}
Potencia máxima disipada en la carga (ecuación ajustada):
\begin{equation}
    P_{L,\text{máx}} = \frac{V_{th}^2}{4 R_{th}} = \frac{(0.07042\text{ V})^2}{4 \cdot (352.465\,\Omega)} = \mathbf{3.517\,\mu\text{W}}
\end{equation}

\vspace{0.4cm}
\begin{tcolorbox}[caja_fase, title=FASE 2: EXTENSIÓN ACTIVA (45 A 90 MIN)]
\textbf{Preamplificación con Fuentes Dependientes y Método $V_{\text{test}} / I_{\text{test}}$}
\end{tcolorbox}

\textbf{Problema de Ingeniería:} \\
La señal de $\qty{70.42}{\milli\volt}$ y potencia de $\qty{3.52}{\micro\watt}$ son insuficientes para excitar un ADC. Conectamos el equivalente de Thévenin a una etapa activa compuesta por una VCVS ($E_1 = 50 V_x$) y una resistencia de salida $R_f$.

\begin{center}
    \begin{circuitikz}[scale=0.9, transform shape]
        % Etapa Sensor
        \draw (0,0) node[ground]{} to[R, l=$R_{th} {=} \qty{352.5}{\ohm}$] (0,3)
              to[V, v=$V_{th} {=} \qty{70.42}{\milli\volt}$] (0,5) -- (2,5);
        \node[above] at (1,5.2) {\textbf{[ EQUIVALENTE SENSOR ]}};
        
        % Etapa Activa
        \draw (2,5) to[short, *-] (3,5) node[above]{Nodo 1};
        \draw (3,5) to[R, l=$R_{in} {=} \qty{1}{\kilo\ohm}$, v=$V_x$] (3,0) node[ground]{};
        \draw (3,5) to[R, l=$R_f {=} \qty{100}{\ohm}$] (7,5) node[above]{Terminal OUT} to[short, *-o] (8,5);
        \draw (7,5) to[cV, v=$E_1 {=} 50V_x$] (7,0) node[ground]{};
        
        \node[above] at (5.5,5.8) {\textbf{[ ETAPA ACTIVA ]}};
    \end{circuitikz}
\end{center}

\subsection*{1. ¿Por qué colapsa el método pasivo tradicional?}
Si un estudiante intenta apagar la fuente dependiente $E_1$ tratándola como un cortocircuito, comete un error físico grave: la magnitud de $E_1$ depende intrínsecamente del voltaje $V_x$. La fuente dependiente inyecta energía dinámica al circuito y modifica la impedancia de salida aparente.

\subsection*{2. Aplicación Rigurosa de la Fuente de Prueba Externa ($V_{\text{test}} = \qty{1.0}{\volt}$)}
Para determinar la verdadera resistencia de Thévenin de la salida ($R_{\text{out}}$):
\begin{enumerate}
    \item Apagamos la fuente \textit{independiente}: $V_{th} = \qty{0}{\volt} \to$ Cortocircuito a Tierra.
    \item Mantenemos la fuente dependiente $E_1 = 50 V_x$ conectada.
    \item Inyectamos $V_{\text{test}} = \qty{1.0}{\volt}$ entre los terminales de salida:
\end{enumerate}

\begin{center}
    \begin{circuitikz}[scale=0.9, transform shape]
        \draw (0,0) node[ground]{} to[R, l=$R_{th} {=} \qty{352.5}{\ohm}$] (0,3) -- (1,3);
        \draw (1,3) to[short, *-] (2,3) node[above]{Nodo 1};
        \draw (2,3) to[R, l=$R_{in} {=} \qty{1}{\kilo\ohm}$, v=$V_x$] (2,0) node[ground]{};
        \draw (2,3) to[R, l=$R_f {=} \qty{100}{\ohm}$] (6,3) node[above]{Nodo C};
        \draw (6,3) to[short, *-o] (7.5,3) node[above]{$V_{\text{test}} {=} \qty{1.0}{\volt}$};
        \draw (6,3) to[cV, v=$E_1 {=} 50V_x$] (6,0) node[ground]{};
    \end{circuitikz}
\end{center}

\textbf{Paso A: Relación de Voltajes en el Nodo 1 ($V_x$)}\\
Las resistencias $R_{th}$ y $R_{\text{in}}$ quedan en paralelo respecto al Nodo 1:
\begin{equation*}
    R_{p1} = R_{th} \parallel R_{\text{in}} = \frac{352.465 \cdot 1000}{352.465 + 1000} = \mathbf{260.61\,\Omega}
\end{equation*}
El circuito queda formado por un divisor de tensión excitado por $V_{\text{test}}$:
\begin{equation*}
    V_x = \qty{1.0}{\volt} \cdot \left( \frac{260.61}{260.61 + 100} \right) = \mathbf{\qty{0.7227}{\volt}}
\end{equation*}

\textbf{Paso B: Tensión y Corrientes Generadas}\\
Corriente $I_1$ que fluye desde $V_{\text{test}}$ hacia el Nodo 1:
\begin{equation*}
    I_1 = \frac{\qty{1.0}{\volt} - \qty{0.7227}{\volt}}{100\,\Omega} = \mathbf{\qty{2.773}{\milli\ampere}}
\end{equation*}

\begin{tcolorbox}[caja_alerta, title=Análisis de Impedancia de Salida]
La fuente dependiente impone: $E_1 = 50(\qty{0.7227}{\volt}) = \qty{36.135}{\volt}$. \\
La corriente total absorbida o entregada por la fuente de prueba demuestra que la impedancia de salida queda gobernada por la etapa activa:
\begin{align}
    I_{\text{test}} &= \frac{V_{\text{test}} - 50 V_x}{R_f} = \frac{1.0 - 36.135}{100} = \mathbf{\qty{-351.35}{\milli\ampere}} \\
    R_{\text{out}} &= \frac{V_{\text{test}}}{I_{\text{test}}} = \frac{\qty{1.0}{\volt}}{\qty{-0.35135}{\ampere}} = \mathbf{-2.846\,\Omega} \nonumber \\ 
    &\quad \rightarrow \text{\textbf{(Resistencia Activa Negativa / Ganancia)}}
\end{align}
\end{tcolorbox}

\newpage
\section*{CÓMPUTO CIENTÍFICO EN PYTHON (NumPy)}
\begin{lstlisting}[language=Python]
import numpy as np
import matplotlib.pyplot as plt

# ==========================================
# FASE 1: SENSOR DE WHEATSTONE (PASIVO)
# ==========================================
Vs = 10.0
R1, R2, R3, R4 = 350.0, 350.0, 350.0, 360.0

# 1. Calculo Thevenin
VA = Vs * (R3 / (R1 + R3))
VB = Vs * (R4 / (R2 + R4))
Vth = VB - VA
Rth = (R1 * R3) / (R1 + R3) + (R2 * R4) / (R2 + R4)

print(f"--- FASE 1: SENSOR PASIVO ---")
print(f"Voltaje Thevenin (Vth) : {Vth*1000:.3f} mV")
print(f"Resistencia Thevenin (Rth): {Rth:.3f} Ohms")

# 2. Barrido de Carga y Curva de Potencia
RL_array = np.linspace(10, 1000, 500)
IL_array = Vth / (Rth + RL_array)
PL_array = (IL_array**2) * RL_array * 1e6  # en microWatts

P_max = (Vth**2) / (4 * Rth) * 1e6
print(f"Potencia Maxima en Carga  : {P_max:.3f} uW en RL = {Rth:.1f} Ohms")

# ==========================================
# FASE 2: ETAPA ACTIVA (Vtest / Itest)
# ==========================================
Rin = 1000.0
Rf = 100.0
k_gain = 50.0
Vtest = 1.0

# Resistencia equivalente en nodo Vx
Rp1 = (Rth * Rin) / (Rth + Rin)
Vx = Vtest * (Rp1 / (Rp1 + Rf))
Itest = (Vtest - k_gain * Vx) / Rf
Rth_activa = Vtest / Itest

print(f"\n--- FASE 2: ETAPA ACTIVA ---")
print(f"Voltaje de control Vx : {Vx:.4f} V")
print(f"Corriente de prueba   : {Itest*1000:.3f} mA")
print(f"Resistencia Activa Rout: {Rth_activa:.3f} Ohms")
\end{lstlisting}

\end{document}
